\documentclass[a4paper,12pt]{article}

\usepackage[T2A]{fontenc}
\usepackage[utf8]{inputenc}
\usepackage[russian]{babel}
\usepackage{amsmath,amssymb,mathtools}
\usepackage{helvet}
\renewcommand{\familydefault}{\sfdefault}
\usepackage{icomma}
\usepackage[a4paper,margin=2cm,headheight=26pt]{geometry}
\usepackage{microtype}
\usepackage{tikz}
\usepackage{tabularx,booktabs}
\usepackage{enumitem}
\usepackage{hyperref}
\usepackage{parskip}
\usepackage{fancyhdr}
\usepackage{setspace}
\usepackage{xcolor}

\definecolor{accent}{HTML}{008080}
\definecolor{darkaccent}{HTML}{004D4D}
\definecolor{textgray}{HTML}{333333}
\definecolor{softgray}{HTML}{F2F5F5}

\hypersetup{
  colorlinks=true,
  linkcolor=darkaccent,
  urlcolor=darkaccent
}

\onehalfspacing
\setlength{\parindent}{0pt}
\setlength{\parskip}{0.55em}
\setlist[itemize]{
  leftmargin=1.6em,
  itemsep=0.25em,
  topsep=0.25em
}
\setlist[enumerate]{
  leftmargin=2em,
  itemsep=0.55em,
  topsep=0.35em
}
\emergencystretch=2em

\pagestyle{fancy}
\fancyhf{}
\fancyhead[L]{\textcolor{textgray}{Экономические задачи: вклады}}
\fancyhead[R]{\textcolor{textgray}{Тимофей Васильченко}}
\fancyfoot[C]{\thepage}
\renewcommand{\headrulewidth}{0.4pt}
\renewcommand{\headrule}{%
  \hbox to\headwidth{%
    \color{accent!55}\leaders\hrule height \headrulewidth\hfill
  }%
}

\newcommand{\key}[1]{\textcolor{darkaccent}{\textbf{#1}}}
\newcommand{\rub}{\,\text{руб.}}
\newcommand{\thrub}{\,\text{тыс. руб.}}

\begin{document}

\begin{titlepage}
  \thispagestyle{empty}
  \centering
  \vspace*{2.2cm}

  {\fontsize{20}{24}\selectfont
    \textcolor{darkaccent}{\textbf{Чек-лист}}\par
  }

  \vspace{0.7cm}

  {\fontsize{25}{30}\selectfont
    \textbf{Банковские вклады и сложные проценты}\par
  }

  \vspace{0.35cm}

  {\Large
    Табличная модель, промежуточные операции, сравнение вкладов\par
  }

  \vfill

  {\large Лёвин Артём Александрович\par}

  \vspace{0.15cm}

  {\large эксклюзивно для Тимофея Васильченко\par}

  \vspace{0.7cm}

  {\large \today\par}

  \vfill

  \begin{tikzpicture}[x=1cm,y=1cm]
    \draw[accent!65,line width=1pt]
      (0,0)
      .. controls (3.2,0.65) and (7.0,-0.55) ..
      (10.2,0)
      .. controls (12.0,0.38) and (14.0,0.38) ..
      (15.7,0);

    \draw[darkaccent!40,line width=0.45pt]
      (0.8,-0.18)
      .. controls (4.0,0.28) and (8.8,-0.35) ..
      (14.9,-0.08);
  \end{tikzpicture}
\end{titlepage}

\section*{1. Базовые обозначения}

Пусть первоначальная сумма вклада равна \(S>0\), годовая ставка
составляет \(P\%\). Вводим обозначения
\[
  p=\frac{P}{100},
  \qquad
  r=1+p=1+\frac{P}{100}.
\]

Число \(r\) называют \key{повышающим коэффициентом}.

\begin{itemize}
  \item Начисленные за год проценты от суммы \(B\):
  \[
    Bp.
  \]

  \item Сумма после начисления процентов:
  \[
    B+Bp=B(1+p)=Br.
  \]

  \item При отсутствии промежуточных операций через \(n\) лет:
  \[
    \boxed{B_n=Sr^n}.
  \]
\end{itemize}

\key{Проверка смысла.}
При положительной ставке \(P>0\) выполняется \(r>1\), поэтому сумма
\(Sr^n\) растёт с увеличением количества лет.

\section*{2. Универсальная таблица по годам}

Если операция совершается после начисления процентов в конце года,
обозначим её через \(d_k\):
\[
  d_k>0
  \quad \text{для пополнения},
  \qquad
  d_k<0
  \quad \text{для снятия}.
\]

Переход от конца \((k-1)\)-го года к концу \(k\)-го года имеет вид
\[
  \boxed{B_k=rB_{k-1}+d_k},
  \qquad
  B_0=S.
\]

\begin{table}[ht]
  \centering
  \caption{Структура расчёта одного года}
  \renewcommand{\arraystretch}{1.35}
  \begin{tabularx}{\linewidth}{@{}cXXX@{}}
    \toprule
    Год
    &
    Сумма в начале года
    &
    Начисленные проценты
    &
    Сумма в конце года
    \\
    \midrule
    \(k\)
    &
    \(B_{k-1}\)
    &
    \(pB_{k-1}\)
    &
    \(B_k=B_{k-1}+pB_{k-1}+d_k=rB_{k-1}+d_k\)
    \\
    \bottomrule
  \end{tabularx}
\end{table}

После последовательного раскрытия рекуррентной формулы получаем
\[
  \boxed{
    B_n
    =
    Sr^n+d_1r^{n-1}+d_2r^{n-2}
    +\dots+d_{n-1}r+d_n
  }.
\]

Чем раньше внесены деньги, тем дольше они участвуют в начислении
процентов. Раннее снятие сильнее уменьшает итоговую сумму.

\section*{3. Алгоритм решения задачи на вклад}

\begin{enumerate}
  \item Выпишите исходную сумму \(S\), ставку \(P\%\), количество лет
  и моменты операций.

  \item Перейдите к коэффициентам
  \[
    p=\frac{P}{100},
    \qquad
    r=1+p.
  \]

  \item Уточните порядок действий внутри каждого года: начисление
  процентов, затем операция по счёту, если именно такой порядок задан
  условием.

  \item Заполните таблицу или примените рекуррентную формулу
  \[
    B_k=rB_{k-1}+d_k.
  \]

  \item Переведите дополнительное условие на математический язык.

  \item Выполните преобразования в общем виде и подставьте числовые
  данные.

  \item Проверьте единицы измерения, знак ответа и требование к целому
  числу.
\end{enumerate}

\section*{4. Снятие и возврат одной суммы}

Пусть вклад открыт на три года. После первого года вкладчик снимает
\(x\), после второго года возвращает \(x\).

Тогда
\[
  B_1=Sr-x,
\]
\[
  B_2=(Sr-x)r+x,
\]
\[
  B_3
  =
  \bigl((Sr-x)r+x\bigr)r
  =
  Sr^3-xr^2+xr.
\]

При отсутствии операций вкладчик получил бы \(Sr^3\). Потеря составляет
\[
  \begin{aligned}
    \Delta S
    &=
    Sr^3-B_3
    \\
    &=
    xr^2-xr
    \\
    &=
    xr(r-1)
    \\
    &=
    xrp.
  \end{aligned}
\]

Для \(P=10\%\), \(r=1{,}1\), \(x=2000\rub\):
\[
  \Delta S
  =
  2000\cdot1{,}1\cdot0{,}1
  =
  \boxed{220\rub}.
\]

\key{Главная идея.}
Возвращённые позже деньги успевают получить меньше процентов, поэтому
одинаковые по величине снятие и пополнение полностью друг друга
не компенсируют.

\section*{5. Сравнение порядка операций}

Рассмотрим два вклада с одинаковыми \(S\), \(r\) и суммой операции
\(x\).

\[
  B_{-+}
  =
  Sr^3-xr^2+xr
  \qquad
  \text{--- сначала снятие, затем пополнение},
\]
\[
  B_{+-}
  =
  Sr^3+xr^2-xr
  \qquad
  \text{--- сначала пополнение, затем снятие}.
\]

Разность итоговых сумм:
\[
  \begin{aligned}
    B_{+-}-B_{-+}
    &=
    2xr^2-2xr
    \\
    &=
    2xr(r-1)
    \\
    &=
    2xrp.
  \end{aligned}
\]

При \(x>0\), \(r>1\) выполняется
\[
  B_{+-}-B_{-+}>0.
\]

Следовательно, больший результат даёт порядок, при котором деньги
сначала вносят и позднее снимают.

При \(x=5000\rub\) и \(P=10\%\):
\[
  B_{+-}-B_{-+}
  =
  2\cdot5000\cdot1{,}1\cdot0{,}1
  =
  \boxed{1100\rub}.
\]

\section*{6. Равные пополнения после первых лет}

Пусть вклад открыт на три года. После начисления процентов в конце
первого и второго годов вносится одна и та же сумма \(a\).
В конце третьего года дополнительного взноса нет.

Получаем
\[
  B_1=Sr+a,
\]
\[
  B_2=(Sr+a)r+a,
\]
\[
  B_3=\bigl((Sr+a)r+a\bigr)r.
\]

После раскрытия скобок:
\[
  \boxed{
    B_3=Sr^3+ar^2+ar
  }.
\]

Удобная сгруппированная форма:
\[
  B_3=Sr^3+a(r^2+r).
\]

Если итоговая сумма известна и равна \(F\), то
\[
  a(r^2+r)=F-Sr^3,
\]
\[
  \boxed{
    a=\frac{F-Sr^3}{r^2+r}
  }.
\]

\subsection*{Пример}

Вклад \(3600\thrub\) открыт под \(10\%\) годовых. После первого
и второго годов вносится одинаковая сумма. Через три года вклад вырос
на \(48{,}5\%\).

Итоговая сумма по условию:
\[
  F
  =
  3600\cdot1{,}485
  =
  5346\thrub.
\]

По формуле накопления
\[
  5346
  =
  3600\cdot1{,}1^3
  +
  a\bigl(1{,}1^2+1{,}1\bigr).
\]

Следовательно,
\[
  a
  =
  \frac{5346-4791{,}6}{1{,}21+1{,}1}
  =
  \frac{554{,}4}{2{,}31}
  =
  \boxed{240\thrub}.
\]

\section*{7. Перевод процентов с русского языка}

Если величина \(S_{\text{стар}}\) изменилась на \(q\%\), то
\[
  S_{\text{нов}}
  =
  S_{\text{стар}}
  \left(
    1\pm\frac{q}{100}
  \right).
\]

При увеличении используется коэффициент
\[
  1+\frac{q}{100}.
\]

При уменьшении используется коэффициент
\[
  1-\frac{q}{100}.
\]

Например, увеличение на \(48{,}5\%\) означает умножение на
\[
  1+\frac{48{,}5}{100}
  =
  1{,}485.
\]

Уменьшение на \(17\%\) означает умножение на
\[
  1-\frac{17}{100}
  =
  0{,}83.
\]

\section*{8. Сравнение вкладов и подбор целой ставки}

Пусть вклад A приносит \(10\%\) ежегодно в течение трёх лет.
Вклад B приносит \(11\%\) в первые два года и \(P\%\) в третий год.
Начальные суммы одинаковы и положительны.

Для вклада A:
\[
  B_A
  =
  S\cdot1{,}1^3
  =
  1{,}331S.
\]

Для вклада B:
\[
  B_B
  =
  S\cdot1{,}11^2
  \left(
    1+\frac{P}{100}
  \right).
\]

Условие большей выгоды вклада B:
\[
  S\cdot1{,}11^2
  \left(
    1+\frac{P}{100}
  \right)
  >
  1{,}331S.
\]

Поскольку \(S>0\), деление обеих частей на \(S\) сохраняет знак
неравенства:
\[
  1{,}2321
  \left(
    1+\frac{P}{100}
  \right)
  >
  1{,}331.
\]

Отсюда
\[
  1+\frac{P}{100}
  >
  \frac{1{,}331}{1{,}2321},
\]
\[
  P
  >
  100
  \left(
    \frac{1{,}331}{1{,}2321}-1
  \right).
\]

Точное пороговое значение:
\[
  P
  >
  \frac{98900}{12321}
  \approx
  8{,}0269.
\]

Наименьшая целая ставка:
\[
  \boxed{9\%}.
\]

\begin{center}
  \begin{tikzpicture}[x=1.05cm,y=1cm]
    \draw[->,line width=0.8pt]
      (0,0)--(10.4,0)
      node[right] {\(P,\%\)};

    \foreach \x/\lab in {
      1/1,
      2/2,
      3/3,
      4/4,
      5/5,
      6/6,
      7/7,
      8/8,
      9/9
    }{
      \draw
        (\x,0.10)--(\x,-0.10)
        node[below=3pt] {\lab};
    }

    \draw[darkaccent,line width=1.4pt]
      (8.03,0)--(10.15,0);

    \draw[
      fill=white,
      draw=darkaccent,
      line width=1pt
    ]
      (8.03,0) circle (2.2pt);

    \fill[accent]
      (9,0) circle (2.5pt);

    \node[above=5pt,text=darkaccent]
      at (8.03,0)
      {\(8{,}0269\ldots\)};
  \end{tikzpicture}
\end{center}

\key{Правило отбора.}
При строгом условии \(P>\alpha\) наименьшее целое значение равно
\[
  \lfloor\alpha\rfloor+1.
\]

Смысл условия определяет направление выбора целого ответа.

\section*{9. Полезные формулы}

\subsection*{Одна операция после \(k\)-го года}

Если после \(k\)-го года внесена сумма \(a\), то к концу \(n\)-го года
она превратится в
\[
  ar^{n-k}.
\]

Если после \(k\)-го года снята сумма \(a\), её влияние на итоговую
сумму равно
\[
  -ar^{n-k}.
\]

\subsection*{Равные взносы после каждого из первых \(n-1\) лет}

Если после каждого из первых \(n-1\) лет вносится сумма \(a\), то
\[
  B_n
  =
  Sr^n
  +
  a\left(
    r^{n-1}+r^{n-2}+\dots+r
  \right).
\]

Сумма степеней является геометрической прогрессией:
\[
  r+r^2+\dots+r^{n-1}
  =
  \frac{r^n-r}{r-1},
  \qquad
  r\ne1.
\]

Следовательно,
\[
  B_n
  =
  Sr^n
  +
  a\frac{r^n-r}{r-1}.
\]

\subsection*{Равные взносы после каждого года, включая последний}

Если сумма \(a\) вносится после каждого из \(n\) начислений процентов,
то
\[
  B_n
  =
  Sr^n
  +
  a\left(
    r^{n-1}+r^{n-2}+\dots+r+1
  \right),
\]
\[
  B_n
  =
  Sr^n
  +
  a\frac{r^n-1}{r-1},
  \qquad
  r\ne1.
\]

Последний взнос входит в итоговую сумму без множителя \(r\), поскольку
после него проценты ещё не начислялись.

\section*{10. Типичные ошибки и контроль}

\begin{itemize}
  \item Путаница между процентной ставкой \(P\%\) и десятичной дробью
  \[
    p=\frac{P}{100}.
  \]

  \item Начисление процентов на сумму, внесённую в конце того же года.

  \item Потеря повышающего коэффициента у раннего взноса или снятия.

  \item Использование формулы \(Sr^n\) при наличии промежуточных
  операций.

  \item Подмена строгого сравнения \(>\) сравнением \(\ge\).

  \item Деление неравенства на переменную без проверки её знака.
  Для суммы вклада фиксируется условие
  \[
    S>0.
  \]

  \item Округление по десятичным знакам без учёта смысла задачи.

  \item Смешение рублей и тысяч рублей.

  \item Добавление операции в год, где условие её не предусматривает.

  \item Преждевременная подстановка чисел, приводящая к громоздким
  вычислениям.

  \item Пропуск проверки ответа подстановкой в исходную модель.
\end{itemize}

\key{Финальная проверка решения}

\begin{enumerate}
  \item Все операции расположены в правильных годах.
  \item Знаки пополнений и снятий выбраны верно.
  \item Каждая операция получила правильную степень коэффициента \(r\).
  \item Деление выполнено только на заведомо ненулевую величину.
  \item При решении неравенства учтён знак делителя.
  \item Целый ответ выбран по смыслу условия.
  \item Денежные единицы указаны в ответе.
\end{enumerate}

\newpage

\section*{Тренировочный блок}

\subsection*{Средний уровень}

\begin{enumerate}
  \item Вклад \(80000\rub\) открыт под \(8\%\) годовых на три года.
  Промежуточных операций нет. Найдите сумму на счёте в конце третьего
  года.

  \item Вклад открыт под \(12\%\) годовых на три года. После начисления
  процентов в конце первого года вкладчик снял \(3000\rub\).
  Других операций не было. На сколько рублей итоговая сумма стала
  меньше суммы по первоначальному плану?

  \item Вкладчик положил \(100000\rub\) под \(10\%\) годовых.
  После начисления процентов в конце первого и второго годов он вносил
  по \(5000\rub\). Найдите сумму на счёте в конце третьего года.
\end{enumerate}

\subsection*{Уровень выше среднего}

\begin{enumerate}[resume]
  \item Два вкладчика положили одинаковые суммы под \(9\%\) годовых
  на три года. Первый после первого года снял \(4000\rub\), после
  второго года внёс \(4000\rub\). Второй выполнил эти операции
  в обратном порядке. Кто получил больше и на сколько рублей?

  \item Владимир положил \(3600\thrub\) под \(10\%\) годовых.
  После начисления процентов в конце первого и второго годов он вносил
  одну и ту же сумму. В конце третьего года вклад оказался на
  \(48{,}5\%\) больше первоначального. Найдите ежегодное пополнение.

  \item Вклад A ежегодно увеличивается на \(8\%\) в течение трёх лет.
  Вклад B увеличивается на \(9\%\) в первые два года и на целое число
  процентов \(P\) в третий год. Начальные суммы одинаковы.
  Найдите наименьшее целое \(P\), при котором вклад B выгоднее.

  \item Вклад открыт под \(10\%\) годовых. После первого и второго
  годов вкладчик вносил по \(20000\rub\). Через три года на счёте
  оказалось \(179300\rub\). Найдите первоначальную сумму.

  \item Вклад открыт под \(10\%\) годовых на три года. После первого
  года вкладчик снял \(x\rub\), после второго года внёс \(x\rub\).
  Итоговая сумма оказалась на \(550\rub\) меньше запланированной.
  Найдите \(x\).

  \item На вклад положили \(200\thrub\). После первого года, сразу
  после начисления процентов, внесли ещё \(22\thrub\).
  После второго года на счёте стало \(266{,}2\thrub\).
  Годовая ставка в оба года была одинаковой. Найдите ставку.
\end{enumerate}

\subsection*{Сложный уровень}

\begin{enumerate}[resume]
  \item На вклад A положили \(200000\rub\) под \(10\%\) годовых
  на три года. На вклад B положили ту же сумму.

  В первый год ставка вклада B равна \(12\%\); после начисления
  процентов сняли \(10000\rub\).

  Во второй год ставка равна \(8\%\); после начисления процентов
  внесли \(10000\rub\).

  В третий год ставка составляет целое число процентов \(P\),
  промежуточных операций нет.

  Найдите наименьшее \(P\), при котором итоговая сумма вклада B
  строго больше итоговой суммы вклада A.
\end{enumerate}

\newpage

\section*{Ответы к упражнениям}

\begin{table}[ht]
  \centering
  \caption{Краткие ответы}
  \renewcommand{\arraystretch}{1.4}
  \begin{tabularx}{\linewidth}{@{}cX@{}}
    \toprule
    № & Ответ \\
    \midrule
    1 & \(100776{,}96\rub\). \\
    2 & \(3763{,}20\rub\). \\
    3 & \(144650\rub\). \\
    4 & Второй вкладчик; на \(784{,}80\rub\). \\
    5 & \(240\thrub\). \\
    6 & \(7\%\). \\
    7 & \(100000\rub\). \\
    8 & \(5000\rub\). \\
    9 & \(10\%\). \\
    10 & \(11\%\). \\
    \bottomrule
  \end{tabularx}
\end{table}

\end{document}